% ===================================================================
%  7. TAULA — Herria neguan. — Baserria udaberrian.
%  Traduction 2026 d'apres texto/io, controlee sur texto/fr, texto/en
%  et texto/es. Consigne : voir texto/eu/10-taula-01.tex.
%
%  OUVERTURE DE LA DEUXIEME SERIE : « BIGARREN SAILA ».
%
%  LE TABLEAU DE LA FERME EST CELUI OU LE BASQUE EST LE PLUS CHEZ LUI
%  DE TOUT LE VOLUME. Le baserri est l'institution meme du pays, et
%  chaque objet du cour de ferme a son nom propre, sans compose et sans
%  emprunt : « baserritarra » le fermier, « artzaina » le berger,
%  « behia » la vache, « ardia » la brebis, « astoa » l'ane, « putzua »
%  le puits, « lastategia » la grange. C'est l'inverse exact du tableau
%  de la marine, qui viendra au 10.
% ===================================================================

%%K t07-noto-1 noto
\VUnotes{143.2mm}{%
(*) Bazkariaren xehetasunak ikus 6. eta 13. tauletan.%
}

%%K t07-apar-1 apar
\VUsaut{4.0mm}
\VUcentre{13.2pt}{90}{BIGARREN SAILA}
\VUsaut{3.0mm}
\VUfilet{16mm}
\VUsaut{5.6mm}
\VUtitre{15.0pt}{60}{7. TAULA}
\VUsaut{3.0mm}

%%K t07-tit-1 sub
\VUcentre{12.0pt}{0}{\textit{Lehen agerraldia.}}
\VUsaut{3.0mm}

%%K t07-tit-2 sub
\VUpk{10.2pt}{40}{Herria neguan.}
\VUsaut{4.0mm}

%%K t07-c1-01-1 p
1. --- Urte berriko jaietan aitarekin joan nintzen Etxaberrira, herrixka batera,
non berak eginkizunak baitzituen.

%%K t07-c1-02-1 p
2. --- \VUgras{Diligentzian} \textsuperscript{(11)} joan ginen, hiriaren eta herri
horren artean ibiltzen denean; baina oso hotz egiten zuenez, bidaia ez zen batere
atsegina izan hain gurdi deserosoan.

%%K t07-c1-03-1 p
3. --- Beraz, benetan poztu ginen \VUgras{kotxezainak} bere diligentzia plazan
gelditu zuenean, \textit{«Hartz Zuria»} \VUgras{ostatuaren} \textsuperscript{(2)}
aurrean. \VUgras{Ostalaria} \textsuperscript{(4)} atarian gure zain zegoen, bere
\VUgras{pipa} lasai erretzen.

%%K t07-c1-03-2 p
Sartzeko gonbita egin zigun, eta ez zidaten bitan eskatu behar izan sutegian
kraskatzen zen mahatsondo-suaren aurrean esertzeko. Orduan bost axola zitzaidan
\VUgras{ipar haize} zorrotza, zeinak \VUgras{iragarki} zaharra
\textsuperscript{(3)} kirrinkarazten baitzuen eta tximiniatik ateratzen zen
\VUgras{kea} \textsuperscript{(1)} moltso beltzetan eramaten.

%%K t07-c1-04-1 p
4. --- Bazkalordua zen. Berandutu gabe, merezi zuen ohore osoa eman genion eman
ziguten janari xume baina goxoari (*).

%%K t07-tit-3 sub
\VUpk{10.2pt}{40}{Bisita batzuk.}
\VUsaut{3.0mm}

%%K t07-c1-05-1 p
5. --- Bazkalondoren aitarekin joan nintzen, aseguru-agentea baita, haren
bezeroengana. Lehenik \VUgras{errementaria} \textsuperscript{(5)} bisitatu genuen,
ostatuaren ondoan bizi dena. Zaldi bat ferratzen ari zen. Harritu ninduen haren
morroiaren indarrak, zeinak zaldiaren \VUgras{apatxa} sendo eusten baitzion bere
larruzko mantalaren kontra, errementariak mailu-kolpe astunez \VUgras{ferra}
\textsuperscript{(6)} iltzatzen zuen bitartean.

%%K t07-c1-06-1 p
6. --- Gero herriko \VUgras{zapatariarengana} \textsuperscript{(7)} joan ginen,
Jakobe zaharrarengana. Nahiz eta iragarkitzat \VUgras{bota} bikain bat zuen,
zapata higatuen pare zahar bat konpontzen aski zuen.

%%K t07-c1-06-2 p
Zaharrak barrez esan zigun: «Hirian „oinetakogilea” nintzen, eta ez nuen
bezerorik. Hemen „konpontzailea” naiz, eta lana soberan dut».

%%K t07-c1-07-1 p
7. --- Bere auzokoaz hitz egin zigun, \VUgras{bizargileaz} \textsuperscript{(8)},
zeinaren bezero guztiak asegabe baitaude. Haren \VUgras{labanak} beti koskatuta
daude, eta haren \VUgras{guraizeek} inoiz ez dute behar bezala mozten. Baina
herriko bizargile bakarra denez, jendeak ez du beste erremediorik haren etxean
bizarra kentzeko eta ilea mozteko. Gainera gizon zekena eta itsusia da.

%%K t07-c1-08-1 p
8. --- Zorionez, beste \VUgras{auzokoak} ez dira haren antzekoak.
\VUgras{Gurdigilea} \textsuperscript{(9)}, adibidez, gizon ona, langilea eta
prestua da. Gurdia konpontzera hari ematea atsegin hutsa da. Haren lana beti dago
txukun eginda.

%%K t07-c1-09-1 p
9. --- Jakobe zaharrak ez zuen \VUgras{larrugileaz} \textsuperscript{(10)} ere
kexarik, zeinari batzuetan \VUgras{uhalak} josten laguntzen baitio, lanak presa
duenean.

%%K t07-c1-09-2 p
Aitak familiaz galdetu zion. «Denak ondo. Nire biloba \VUgras{eskolan}
\textsuperscript{(12)} dago. \VUgras{Andereñoa} \textsuperscript{(13)} oso pozik
dago berarekin; \VUgras{ikasle} \textsuperscript{(14)} ona da. Ez nire seme
gaizto hori bezala; jokoa besterik ez du buruan. \VUgras{Maisuak}
\textsuperscript{(15)} ezin dio ezer ikasarazi».

%%K t07-tit-4 sub
\VUsaut{3.0mm}
\VUpk{10.2pt}{40}{Herriko esamesak.}
\VUsaut{3.0mm}

%%K t07-c1-10-1 p
10. --- Gero zaharrak herriko berriak kontatu zizkigun. Eskola berri bat eraikitzeko
asmoa zuten, \VUgras{udaletxekoa} txikiegia geratu baitzen. Ez zen zaila izango,
errotariaren baratzearen zati bat erostea aski baitzen. Gizajoari oso ondo etorriko
zitzaion. Hotzak hasi zirenetik \VUgras{errotaren} \textsuperscript{(16)}
\VUgras{errotarriek} ezin izan dute alerik ehotu, \VUgras{gurpila}
\textsuperscript{(17)} \VUgras{izotzak} \textsuperscript{(25)} hartu baitu
\VUgras{errekan} \textsuperscript{(23)}.

%%K t07-c1-11-1 p
11. --- «Eraikuntzez ari garela: ikusi duzue gure \VUgras{eliza} berria
\textsuperscript{(18)}? Oso pintoreskoa dago elur-txanoaren azpian.
\VUgras{Beleak} \textsuperscript{(19)}, beren kanpandorre zaharra ezagutzen ez
dutenak, goibel pausatzen dira \VUgras{hilerriko} \textsuperscript{(20)} zuhaitz
handian».

%%K t07-c1-12-1 p
12. --- Hilerriaren aipamenak zapatariari aitaren ezagunak gogorarazi zizkion,
iaz hil zirenak. «Zenbat \VUgras{hilobi} \textsuperscript{(21)}, jauna, gehitu
diren lehengoei \VUgras{altzifrearen} \textsuperscript{(22)} azpian! Heriotzak
gure notario zaharra eraman du. Herri osoa izan zen haren \VUgras{ehorzketan}».

%%K t07-c1-13-1 p
13. --- «Hara haren oinordekoa errekan irristaka. Lehengo egunean etorri zen bere
\VUgras{irristailuaren} \textsuperscript{(26)} uhala konpon niezaion, hautsi
baitzitzaion».

%%K t07-c1-14-1 p
14. --- «Eta hor, errekaren ertzean, \VUgras{herri-zaindari} berria
\textsuperscript{(27)}. Jendarme ohia da eta herriko mutil bihurrien izua. Ihes
egiten dute haren \VUgras{polainak} \textsuperscript{(29)} eta haren
\VUgras{txanoa} \textsuperscript{(28)} ikusi orduko».

%%K t07-c1-15-1 p
15. --- «A! hor Joseba zaharra \VUgras{txotx-gurdi} bat \textsuperscript{(30)}
daramala okinari. --- Egia da, esan zuen aitak, ezagutzen dut haren \VUgras{mando}
uztarketa \textsuperscript{(31)}. Berarekin hitz egin behar dut, eta hau
baliatuko dut. Agur, Jakobe zaharra».

%%K t07-c1-16-1 p
16. --- Ateratzen ari ginenean, herriko haurrak eskolatik irteten ari ziren eta
\VUgras{izotz-bidexkara} \textsuperscript{(33)} jauzten, erori eta
\VUgras{erorikoan} \textsuperscript{(34)} zerbait hausteko arriskuan. Haietako
batek gurdigilearengandik bere \VUgras{lera} \textsuperscript{(32)} hartu zuen,
lagun bat gainean jarri eta korrika hasi zen.

%%K t07-c1-17-1 p
17. --- Beste batzuk \VUgras{elur-metara} \textsuperscript{(37)} oldartu
ziren, \VUgras{zubiaren} \textsuperscript{(40)} ondokora, eta \VUgras{elur-panpina}
\textsuperscript{(39)} bonbardatzen hasi, bezperan egin zutena eta zeinari kapela,
pipa eta eskola-zorroa eman baitzizkioten.

%%K t07-c1-18-1 p
18. --- Bost axola zitzaien haize izoztua eta hotz gorria. Gehienek ez zuten
\VUgras{bufandarik} \textsuperscript{(35)} ere, eta \VUgras{elur-bola}
\textsuperscript{(38)} jokoaren ondoren berotzeko, aski zuten \VUgras{gaztaina}
\textsuperscript{(42)} erre eskukada bat, Jakelina zaharrari erosia,
\VUgras{saltzaileari}, zeina hotzak dardarka baitago bere \VUgras{xal}
meheegiaren \textsuperscript{(43)} azpian.

%%K t07-tit-5 sub
\VUsaut{3.0mm}
\VUcentre{12.0pt}{0}{\textit{Bigarren agerraldia.}}
\VUsaut{3.0mm}

%%K t07-tit-6 sub
\VUpk{10.2pt}{40}{Baserria udaberrian.}
\VUsaut{4.0mm}

%%K t07-c2-01-1 p
1. --- Neguaren poz guztiekin ere, poz sakonez hartzen dugu \VUgras{udaberriaren}
itzulera.

%%K t07-c2-01-2 p
Zeinen koadro atsegina den baserriko patioa arratsean, maiatzean!

%%K t07-c2-02-1 p
2. --- Zeruertzean \VUgras{eguzkia} \textsuperscript{(1)} astiro sartzen ari da,
\VUgras{landak} \textsuperscript{(3)} \VUgras{izpi} gorriz \textsuperscript{(2)}
bustiz. \VUgras{Nekazari} langilea \textsuperscript{(6)} presaka ari da bere
\VUgras{soroa} \textsuperscript{(4)} goldatzen.

%%K t07-c2-02-2 p
Bere \VUgras{goldeaz} \textsuperscript{(5)}, zeina \VUgras{zaldi} sendo batek
\textsuperscript{(7)} tiratzen baitu, \VUgras{ildoa} \textsuperscript{(8)} egiten
du, zeinetan \VUgras{ereileak} \textsuperscript{(9)} eskukadaka barreiatuko baitu
\VUgras{hazia}, urrezko galburuak emango dituena.

%%K t07-c2-03-1 p
3. --- \VUgras{Segalariek} \textsuperscript{(11)} beren segekin belardiko belarra
ebaki dute, itzultzaileek \VUgras{belarra} \textsuperscript{(10)} lehortu dute,
\VUgras{sardez} \textsuperscript{(12)} eta eskuarez irauliz, eta hara non
etxeratzen diren.

%%K t07-c2-03-2 p
Batzuk idiek tiratutako gurdi astunaren aurretik doaz, tresnak bizkarrean. Beste
batzuk, alferragoak, eroso jarri dira belar usaintsuaren gainean.

%%K t07-c2-04-1 p
4. --- Igarotzean, adiskidetsu agurtzen dute \VUgras{lorezaina}
\textsuperscript{(16)}, zeina \VUgras{baratzean} \textsuperscript{(13)} ari baita
lanean, \VUgras{fruta-arbolak} inausten eta \VUgras{udareondoak}
\textsuperscript{(14)} txertatzen. \VUgras{Aranondoak} \textsuperscript{(15)}
loretan daude, eta ondo ongarritutako \VUgras{ortuan} \textsuperscript{(17)}
barazkiak, aza eta uraza ederrak daude jada.

%%K t07-c2-04-2 p
Hesi baxuaren gainean \VUgras{pauma} eder batek \textsuperscript{(18)} isatsa
zabaltzen du, bere luma bikainak erakusteko.

%%K t07-tit-7 sub
\VUpk{10.2pt}{40}{Baserriko patioa.}
\VUsaut{3.0mm}

%%K t07-c2-05-1 p
5. --- \VUgras{Ukuiluaren} \textsuperscript{(19)} ateak zabalik daude, eta
askak eta \VUgras{ohantzea} \textsuperscript{(23)} ikusten dira,
\VUgras{behorrarenak} \textsuperscript{(21)}, zeina \VUgras{zaldizainak}
\textsuperscript{(20)} askatu berri baitu. \VUgras{Moxala}
\textsuperscript{(22)}, hanka luzeetan hain barregarria, amaren atzetik doa,
jauzika.

%%K t07-c2-06-1 p
6. --- \VUgras{Putzuaren} \textsuperscript{(29)} ondoan \VUgras{behiak}
\textsuperscript{(25)} edatera doaz zurezko \VUgras{aska} batetik
\textsuperscript{(26)}, zeina edaleku baitzaie. \VUgras{Txahalak}
\textsuperscript{(24)} aukera baliatzen du edoskitzeko, eta \VUgras{idiak}
\textsuperscript{(27)}, bazkaren usain ona sumaturik, alai marrumatzen du eta
harro altxatzen du burua bere \VUgras{adar} indartsuekin \textsuperscript{(28)}.

%%K t07-c2-07-1 p
7. --- Denak lanean daude. \VUgras{Langile emakumeak} \textsuperscript{(32)}
putzuko \VUgras{soka} \textsuperscript{(31)} eusten du. Ura ateratzen du
\VUgras{ontziaz} \textsuperscript{(30)}, abereei edaten emateko.
\VUgras{Lastategiaren} \textsuperscript{(33)} ondoan gizon batek barreiatutako
lastoa biltzen du, eta \VUgras{baserriaren} \textsuperscript{(34)} atearen
aurrean \VUgras{iruleak} \textsuperscript{(35)} bere gurpilarekin eta
\VUgras{kilorearekin} \VUgras{lihoa} edo \VUgras{kalamua} irauten du, antzinako
egunetan bezala.

%%K t07-tit-8 sub
\VUsaut{3.0mm}
\VUpk{10.2pt}{40}{Baserritarra.}
\VUsaut{3.0mm}

%%K t07-c2-08-1 p
8. --- \VUgras{Baserritarrak} \textsuperscript{(36)} lan hori guztia zuzentzen du
eta bere jendeari aginduak ematen dizkio. \VUgras{Area}
\textsuperscript{(40)} eta \VUgras{aitzurra} \textsuperscript{(39)} teilapean
gordetzeko agintzen du: \VUgras{etxolaren} \textsuperscript{(38)} ondoan utzi
dituzte, \VUgras{txakurrarenaren} \textsuperscript{(37)} ondoan.

%%K t07-c2-09-1 p
9. --- Haren oihuek \VUgras{usoa} \textsuperscript{(41)} izutzen dute, zeina
lasterka baitoa \VUgras{usategira} \textsuperscript{(42)}, eta \VUgras{oiloak}
\textsuperscript{(43)}, zeinak presaka baitoaz \VUgras{oilategira}
\textsuperscript{(44)}.

%%K t07-c2-10-1 p
10. --- \VUgras{Artegiaren} \textsuperscript{(48)} aurrean \VUgras{artzain}
zaharrak \textsuperscript{(45)} bere \VUgras{artaldea} \textsuperscript{(49)}
biltzen du. Bere \VUgras{makila} \textsuperscript{(46)} eskuan, zeina inoiz ez
baita berarengandik banatzen, ezta bere \VUgras{blusa} zurbila
\textsuperscript{(47)} ere, eskortara eramaten ditu \VUgras{ahariak}
\textsuperscript{(50)}, \VUgras{ardiak} \textsuperscript{(53)},
\VUgras{bildotsak} \textsuperscript{(54)}, \VUgras{ahuntzak}
\textsuperscript{(51)} eta \VUgras{antxumeak} \textsuperscript{(52)}. Bere
\VUgras{txakur} leiala \textsuperscript{(55)}, Argus, azkena doa eta zaunkaz
bizkortzen ditu atzeratuak.

%%K t07-c2-11-1 p
11. --- Zarata eta iskanbila horrek guztiak ez du asaldatzen \VUgras{esne-behi}
onaren \textsuperscript{(56)} lasaitasuna, zeinak bere \VUgras{joarea}
\textsuperscript{(57)} dilindatzen baitu \VUgras{behitegiaren}
\textsuperscript{(58)} atearen aurrean jezten duten bitartean.

%%K t07-c2-12-1 p
12. --- Badirudi ulertzen duela esnea eman behar duela, \VUgras{baserritar
emakumeak} \textsuperscript{(60)} \VUgras{gurina} \VUgras{gurin-ontzian}
\textsuperscript{(61)} irabiatu dezan edo \VUgras{gazta} bikainak
\textsuperscript{(64)} egin ditzan, zeinak \VUgras{upelaren}
\textsuperscript{(63)} gainean zeharka jarritako ohol batean lehortzen baitira.

%%K t07-c2-13-1 p
13. --- \VUgras{Esnetegi} osoa \textsuperscript{(59)} garbi mantentzen du
\VUgras{morroiak} \textsuperscript{(65)}, zeinak \VUgras{esne-ontziak}
\textsuperscript{(62)} garbitu berri baititu daraman ontzi handian.

%%K t07-tit-9 sub
\VUsaut{3.0mm}
\VUpk{10.2pt}{40}{Hegazti-korta.}
\VUsaut{3.0mm}

%%K t07-c2-14-1 p
14. --- Bere zurezko oinetako astunekin nahiko baldar dabil, eta horregatik
izutzen ditu \VUgras{indioilarra} \textsuperscript{(68)}, \VUgras{ahateak}
\textsuperscript{(66)}, \VUgras{ahate arrak} \textsuperscript{(67)} eta
\VUgras{antzarak} \textsuperscript{(69)}, zeinek zabal irekitzen baitituzte beren
\VUgras{mokoak} \textsuperscript{(70)} eta zalapartaka hasten.

%%K t07-c2-14-2 p
\VUgras{Oilarrak} \textsuperscript{(71)}, gerlariagoa, bere \VUgras{gandor}
gorria \textsuperscript{(72)} altxatzen du, \VUgras{isatseko}
\textsuperscript{(73)} lumak astintzen ditu eta «kukurruku» ozena jotzen.

%%K t07-c2-15-1 p
15. --- \VUgras{Oilo lokak} \textsuperscript{(74)} \VUgras{simaur-pilan}
\textsuperscript{(81)} arakatzen du bere \VUgras{txitekin}
\textsuperscript{(75)}. \VUgras{Txerrikumea} \textsuperscript{(78)} amarengana
doa lasterka, \VUgras{txerri} handi harengana \textsuperscript{(77)}, zeina
ukuiluaren ondoan ikusten baita.

%%K t07-c2-16-1 p
16. --- Beren kaioletan \VUgras{untxiek} \textsuperscript{(79)} gogotik jaten
dute \VUgras{belar} samurra \textsuperscript{(80)}, baserritarraren alabak
ekartzen diena. Eugeniak oso maite ditu bere untxiak eta egunero, oilategian
\VUgras{arrautza} freskoak \textsuperscript{(76)} hartu ondoren, bere adiskide
txikiak bisitatzera etortzen da.
\VUsaut{6.0mm}
\VUfilet{22mm}
